% !TEX program = xelatex
% Đặc tả chi tiết hệ thống — Social Media Real-Time System (SE347)
% Biên dịch bằng XeLaTeX (chạy 2 lần liên tiếp để mục lục hiển thị đúng):
%     xelatex dac-ta-he-thong.tex
%     xelatex dac-ta-he-thong.tex
% Nếu máy có cài font "Times New Roman" (thường có sẵn trên Windows) và muốn
% dùng đúng font này thay vì font thay thế bên dưới, đổi dòng \setmainfont.
\documentclass[12pt,a4paper]{report}

\usepackage{fontspec}
\usepackage{polyglossia}
\setmainlanguage{vietnamese}
\setmainfont{TeX Gyre Termes} % thay bằng {Times New Roman} nếu máy có sẵn font này

\usepackage[a4paper,left=3cm,right=2cm,top=2.5cm,bottom=2.5cm]{geometry}
\usepackage{setspace}
\onehalfspacing

\usepackage{longtable}
\usepackage{booktabs}
\usepackage{array}
\usepackage{enumitem}
\usepackage{xcolor}
\usepackage{url}
\urlstyle{tt}
\usepackage{fancyhdr}
\usepackage[hidelinks]{hyperref}

\hypersetup{
    pdftitle={Đặc tả chi tiết hệ thống Social Media Real-Time System},
    pdfauthor={Sinh viên thực hiện},
    pdfsubject={Đồ án 1 - SE347},
    bookmarksnumbered=true,
    bookmarksopen=true
}

\setlength{\parindent}{0pt}
\setlength{\parskip}{6pt}

\pagestyle{fancy}
\fancyhf{}
\fancyhead[L]{\slshape\leftmark}
\fancyhead[R]{\thepage}
\renewcommand{\headrulewidth}{0.4pt}
\renewcommand{\footrulewidth}{0pt}

\fancypagestyle{plain}{%
  \fancyhf{}
  \fancyhead[L]{\slshape\leftmark}
  \fancyhead[R]{\thepage}
  \renewcommand{\headrulewidth}{0.4pt}
}

% Bảng thực thể / API dùng chung một khuôn dạng longtable 3 cột.
\newcommand{\entityTableHeader}[1]{%
    \toprule
    \textbf{Trường} & \textbf{Kiểu dữ liệu} & \textbf{Mô tả} \\
    \midrule
    \endfirsthead
    \multicolumn{3}{l}{\small\itshape #1 (tiếp theo)} \\
    \toprule
    \textbf{Trường} & \textbf{Kiểu dữ liệu} & \textbf{Mô tả} \\
    \midrule
    \endhead
    \bottomrule
    \endfoot
}

\newcommand{\apiTableHeader}[1]{%
    \toprule
    \textbf{Method} & \textbf{Endpoint (qua Gateway)} & \textbf{Mô tả} \\
    \midrule
    \endfirsthead
    \multicolumn{3}{l}{\small\itshape #1 (tiếp theo)} \\
    \toprule
    \textbf{Method} & \textbf{Endpoint (qua Gateway)} & \textbf{Mô tả} \\
    \midrule
    \endhead
    \bottomrule
    \endfoot
}

% Bảng 2 cột dùng chung (công nghệ, biến môi trường, cổng dịch vụ, ...)
\newcommand{\twocolTableHeader}[3]{%
    \toprule
    \textbf{#2} & \textbf{#3} \\
    \midrule
    \endfirsthead
    \multicolumn{2}{l}{\small\itshape #1 (tiếp theo)} \\
    \toprule
    \textbf{#2} & \textbf{#3} \\
    \midrule
    \endhead
    \bottomrule
    \endfoot
}

% Bảng 3 cột dùng chung với tiêu đề cột tuỳ biến (cổng/DB, biến môi trường...)
\newcommand{\threecolTableHeader}[4]{%
    \toprule
    \textbf{#2} & \textbf{#3} & \textbf{#4} \\
    \midrule
    \endfirsthead
    \multicolumn{3}{l}{\small\itshape #1 (tiếp theo)} \\
    \toprule
    \textbf{#2} & \textbf{#3} & \textbf{#4} \\
    \midrule
    \endhead
    \bottomrule
    \endfoot
}

% Bảng 4 cột dùng chung (dùng cho bảng định tuyến API Gateway)
\newcommand{\fourcolTableHeader}[5]{%
    \toprule
    \textbf{#2} & \textbf{#3} & \textbf{#4} & \textbf{#5} \\
    \midrule
    \endfirsthead
    \multicolumn{4}{l}{\small\itshape #1 (tiếp theo)} \\
    \toprule
    \textbf{#2} & \textbf{#3} & \textbf{#4} & \textbf{#5} \\
    \midrule
    \endhead
    \bottomrule
    \endfoot
}

\title{\Huge\textbf{Đặc tả chi tiết hệ thống}}
\author{}
\date{}

\begin{document}

% ============================================================
% TRANG BÌA
% ============================================================
\begin{titlepage}
\centering
{\large ĐẠI HỌC QUỐC GIA TP. HỒ CHÍ MINH \par}
{\large TRƯỜNG ĐẠI HỌC CÔNG NGHỆ THÔNG TIN \par}
\vspace{0.3cm}
{\large \textbf{KHOA CÔNG NGHỆ PHẦN MỀM} \par} % TODO: sửa lại đúng khoa/bộ môn nếu khác

\vspace{2cm}
\rule{0.8\textwidth}{0.4pt}\par
\vspace{0.4cm}
{\Huge\bfseries ĐẶC TẢ CHI TIẾT HỆ THỐNG\par}
\vspace{0.4cm}
{\LARGE\bfseries MẠNG XÃ HỘI THỜI GIAN THỰC\par}
{\Large (Social Media Real-Time System)\par}
\vspace{0.4cm}
\rule{0.8\textwidth}{0.4pt}\par

\vspace{1.2cm}
{\large Kiến trúc Microservices trên nền Java Spring Boot / Spring Cloud\par}

\vspace{2cm}
{\large Học phần: \textbf{SE347} \par} % TODO: điền đúng tên học phần nếu khác
{\large Đồ án 1 \par}

\vspace{2cm}
\begin{tabular}{ll}
\textbf{Giảng viên hướng dẫn:} & \underline{\hspace{6cm}} \\[0.6cm]
\textbf{Sinh viên thực hiện:} & \underline{\hspace{6cm}} \\[0.3cm]
\textbf{MSSV:} & 22521253 \\[0.6cm]
\textbf{Lớp:} & \underline{\hspace{6cm}} \\
\end{tabular}

\vfill
{\large TP. Hồ Chí Minh, \the\year\par}
\end{titlepage}

\tableofcontents
\clearpage

% ============================================================
\chapter{Tổng quan đồ án}
% ============================================================

\section{Bối cảnh và mục tiêu}

Đồ án xây dựng phần \textbf{backend} cho một mạng xã hội theo mô hình Facebook,
bao gồm đầy đủ các nghiệp vụ cốt lõi: đăng bài (post), story, reels, bình luận,
cảm xúc (reaction), kết bạn/theo dõi, nhóm (group), trang (fanpage), hẹn hò
(dating) và nhắn tin thời gian thực (real-time chat). Hệ thống được thiết kế
theo kiến trúc \textbf{microservices}, mỗi nghiệp vụ là một service độc lập,
giao tiếp với nhau qua REST (đồng bộ) và Apache Kafka (bất đồng bộ), có thể
triển khai và mở rộng độc lập.

Mục tiêu chính của đồ án:
\begin{itemize}[nosep]
    \item Thiết kế và hiện thực một hệ backend nhiều service theo đúng các
    nguyên lý microservices: service discovery, cấu hình tập trung, API
    gateway, giao tiếp bất đồng bộ, chịu lỗi (resilience), quan sát hệ thống
    (observability).
    \item Áp dụng xác thực/phân quyền tập trung bằng JWT tại API Gateway.
    \item Xây dựng một giao diện web (Flutter Web) sử dụng trực tiếp các API
    thật để minh chứng luồng nghiệp vụ đầu-cuối.
    \item Thiết lập observability đầy đủ (distributed tracing, log tập trung,
    metrics) và pipeline CI/CD tự động.
\end{itemize}

\section{Phạm vi chức năng}

Hệ thống gồm \textbf{20 module Maven} (1 thư viện dùng chung, 3 hạ tầng
Spring Cloud, 16 business service nghiệp vụ) cùng một frontend Flutter Web độc
lập. Danh sách chức năng nghiệp vụ chính:

\begin{itemize}[nosep]
    \item Đăng ký/đăng nhập, quản lý phiên (access token + refresh token).
    \item Hồ sơ người dùng, kết bạn, theo dõi (follow), chặn (block), gợi ý kết bạn.
    \item Upload/lưu trữ media (ảnh, video) qua đối tượng lưu trữ tương thích S3.
    \item Đăng bài viết (kèm chia sẻ/repost, gắn thẻ người dùng, quyền riêng tư tuỳ biến).
    \item Bình luận (kèm trả lời lồng nhau), cảm xúc (reaction), lưu bài viết (save).
    \item Story (tự xoá sau 24 giờ), Reels (video ngắn).
    \item Nhóm (group) công khai/riêng tư, trang (fanpage).
    \item Hẹn hò (dating): vuốt (swipe), ghép đôi (match) theo thuật toán chấm điểm.
    \item Nhắn tin thời gian thực (WebSocket/STOMP), trạng thái online (presence).
    \item Thông báo thời gian thực (WebSocket/STOMP).
    \item Bảng tin (feed) được xếp hạng theo mức độ tương tác.
    \item Kiểm duyệt nội dung (lọc từ ngữ thô tục, hàng đợi báo cáo vi phạm cho quản trị viên).
    \item Tìm kiếm tổng hợp trên nhiều loại đối tượng (người dùng/nhóm/trang/bài viết).
\end{itemize}

\section{Công nghệ sử dụng}

\begin{longtable}{p{4.2cm}p{10.8cm}}
\twocolTableHeader{Bảng công nghệ sử dụng}{Nhóm}{Công nghệ / thư viện}
Ngôn ngữ \& nền tảng & Java 17, Maven đa module (reactor 20 module) \\
Framework ứng dụng & Spring Boot 3.3.4 \\
Điều phối microservice & Spring Cloud 2023.0.3 (Netflix Eureka, Spring Cloud Gateway, OpenFeign) \\
Web layer & Spring Web MVC (16 business service) + Spring WebFlux (api-gateway, do Spring Cloud Gateway yêu cầu reactive) \\
Truy xuất dữ liệu & Spring Data JPA (PostgreSQL), Spring Data MongoDB, Spring Data Redis \\
Giao tiếp bất đồng bộ & Spring Kafka (Apache Kafka, chế độ KRaft — không cần Zookeeper) \\
Real-time & Spring WebSocket (giao thức STOMP qua SockJS/raw WebSocket) \\
Bảo mật & JWT (thư viện jjwt 0.12.6) — phát hành tại auth-service, xác thực tập trung tại api-gateway \\
Chịu lỗi & Resilience4j (circuit breaker cho các Feign client liên service) \\
Lưu trữ đối tượng & MinIO (tương thích S3) cho media-service \\
Tài liệu API & springdoc-openapi (OpenAPI 3 / Swagger UI, gộp qua gateway) \\
Kiểm thử & JUnit 5, Mockito, Testcontainers (integration test với DB/Kafka thật) \\
Đóng gói \& triển khai & Docker, Docker Compose (30 container: hạ tầng + service + frontend) \\
CI/CD & GitHub Actions — build/test toàn reactor, build \& push image lên GHCR \\
Observability & Zipkin (tracing), Loki + Promtail + Grafana (log tập trung), Prometheus (metrics) \\
Frontend & Flutter Web (biên dịch tĩnh, phục vụ qua Nginx trong container) \\
\end{longtable}

% ============================================================
\chapter{Kiến trúc hệ thống}
% ============================================================

\section{Kiến trúc tổng quan}

Hệ thống theo mô hình microservices cổ điển: mọi request từ client (web/mobile)
đều đi qua một \textbf{API Gateway} duy nhất, gateway xác thực JWT rồi định
tuyến (route) tới đúng business service thông qua \textbf{service discovery}
(Eureka). Các service không gọi thẳng vào nhau qua địa chỉ IP/port cố định mà
tra cứu lẫn nhau qua tên logic đã đăng ký trên Eureka (load-balanced qua
Spring Cloud LoadBalancer, giao thức \path{lb://TEN-SERVICE}).

\begin{longtable}{p{4cm}p{11cm}}
\twocolTableHeader{Sơ đồ các lớp kiến trúc}{Lớp}{Thành phần \& vai trò}
Client & Web (Flutter Web, cổng 3001) — gọi thẳng API Gateway; kiến trúc REST
+ WebSocket cho phép mở rộng thêm client mobile sau này mà không đổi backend. \\
API Gateway & Cổng 8080 — Spring Cloud Gateway (WebFlux). Xác thực JWT, giới
hạn tần suất (rate limiting), định tuyến theo path, gộp Swagger UI của 16
service vào một điểm truy cập duy nhất. \\
Service Discovery & Eureka Server, cổng 8761 — mọi service (kể cả gateway)
đăng ký và tra cứu lẫn nhau tại đây; Prometheus cũng dùng Eureka để tự phát
hiện target cần scrape metrics. \\
Config Server & Cổng 8888 — cấp cấu hình dùng chung (profile \path{native},
đọc từ thư mục \path{config-repo}) cho mọi service có khai báo
\path{spring.config.import}. \\
Business services & 16 service nghiệp vụ độc lập (xem Chương~\ref{ch:services}), mỗi service
một cơ sở dữ liệu riêng (\emph{database-per-service}), giao tiếp ngang hàng
qua OpenFeign (đồng bộ) hoặc Kafka (bất đồng bộ). \\
Message broker & Apache Kafka (chế độ KRaft, không cần Zookeeper) — kênh giao
tiếp bất đồng bộ giữa các service (xem mục~\ref{sec:kafka}). \\
Lưu trữ dữ liệu & PostgreSQL (10 database riêng biệt), MongoDB (4 database),
Redis (feed fanout, presence, rate limiter), MinIO (media nhị phân). \\
Observability & Zipkin, Prometheus, Loki + Promtail + Grafana (xem
Chương~\ref{ch:observability}). \\
\end{longtable}

\section{API Gateway}

API Gateway (module \path{infra/api-gateway}, cổng 8080) là điểm vào duy nhất
của toàn hệ thống, xây dựng trên Spring Cloud Gateway (nền WebFlux — reactive,
khác với các business service dùng Spring MVC truyền thống). Gateway đảm nhận
ba trách nhiệm chính:

\begin{enumerate}[nosep]
    \item \textbf{Xác thực JWT tập trung} — mọi service nghiệp vụ phía sau
    không cần tự giải mã/validate JWT, chỉ cần đọc header \path{X-User-Id} /
    \path{X-User-Roles} do gateway gắn sẵn (xem Chương~\ref{ch:security}).
    \item \textbf{Định tuyến (routing)} theo tiền tố đường dẫn, ánh xạ 1--1
    sang tên service đã đăng ký trên Eureka.
    \item \textbf{Giới hạn tần suất (rate limiting)} theo từng route bằng
    thuật toán token bucket lưu trạng thái trên Redis
    (\path{RequestRateLimiter} filter).
\end{enumerate}

Bảng định tuyến đầy đủ (trích từ \path{application.yml} của api-gateway):

\begin{longtable}{p{3.2cm}p{4cm}p{2.3cm}p{4.5cm}}
\fourcolTableHeader{Bảng định tuyến API Gateway}{Route}{Predicate (path)}{Rate limit mặc định}{Ghi chú}
auth-service & \path{/api/auth/**} & 1 req/s, burst 5 & Giới hạn chặt hơn hẳn để chống brute-force đăng nhập/đăng ký. \\
user-service & \path{/api/users/**} & 20 req/s, burst 40 & \\
media-service & \path{/api/media/**} & 20 req/s, burst 40 & \\
post-service & \path{/api/posts/**} & 20 req/s, burst 40 & \\
comment-service & \path{/api/comments/**} & 20 req/s, burst 40 & \\
reaction-service & \path{/api/reactions/**}, \path{/api/saved/**} & 20 req/s, burst 40 & \\
story-service & \path{/api/stories/**} & 20 req/s, burst 40 & \\
reels-service & \path{/api/reels/**} & 20 req/s, burst 40 & \\
group-service & \path{/api/groups/**} & 20 req/s, burst 40 & \\
fanpage-service & \path{/api/pages/**} & 20 req/s, burst 40 & \\
dating-service & \path{/api/dating/**} & 20 req/s, burst 40 & \\
chat-service & \path{/api/chat/**} & 20 req/s, burst 40 & \\
chat-service (WS) & \path{/ws/**} & 5 req/s, burst 10 & Route riêng, \path{uri: lb:ws://CHAT-SERVICE}. \\
notification-service & \path{/api/notifications/**} & 20 req/s, burst 40 & \\
notification-service (WS) & \path{/ws-notifications/**} & 5 req/s, burst 10 & Route riêng cho WebSocket. \\
feed-service & \path{/api/feed/**} & 20 req/s, burst 40 & \\
moderation-service & \path{/api/moderation/**} & 20 req/s, burst 40 & \\
search-service & \path{/api/search/**} & 20 req/s, burst 40 & \\
\end{longtable}

Ghi chú: mọi giá trị rate limit đều cấu hình được qua biến môi trường
(\path{RATE_LIMIT_REPLENISH}, \path{RATE_LIMIT_BURST}, ...), không hard-code
trong image — xem Phụ lục~\ref{ap:env}. Ngoài các route nghiệp vụ trên, gateway
còn có 15 route phụ dạng \path{/docs/<service>/v3/api-docs} (không giới hạn
tần suất, không yêu cầu xác thực) để gộp tài liệu OpenAPI của từng service vào
một Swagger UI chung tại \path{/swagger-ui.html}.

\section{Service Discovery (Eureka)}

Mọi service (trừ Eureka Server chính nó) khai báo \path{eureka.client.
service-url.defaultZone} trỏ về Eureka Server (cổng 8761) và tự đăng ký khi
khởi động, gỡ đăng ký khi tắt (heartbeat định kỳ). Gateway và các Feign client
tra cứu instance khả dụng qua tên logic (ví dụ \path{lb://POST-SERVICE}) thay
vì địa chỉ cố định — cho phép chạy nhiều instance của cùng một service để cân
bằng tải mà không cần đổi cấu hình phía gọi.

\section{Config Server}

Config Server (module \path{infra/config-server}, cổng 8888) chạy ở profile
\path{native}, đọc cấu hình dùng chung từ thư mục \path{config-repo/
application.yml} (mount vào container qua volume) thay vì cần một Git
repository riêng. Các giá trị ở đây (mức log, danh sách actuator endpoint công
khai...) được merge \emph{dưới} cấu hình cục bộ của từng service — cho phép
đổi các cấu hình xuyên suốt hệ thống (ví dụ hạ mức log) mà không cần build lại
image, chỉ cần sửa file và khởi động lại service tương ứng.

\section{Giao tiếp giữa các service}

\subsection{Giao tiếp đồng bộ — OpenFeign}

Khi một service cần dữ liệu "tươi" ngay lập tức từ service khác (ví dụ
post-service cần kiểm tra quan hệ bạn bè khi enforce quyền riêng tư
\path{FRIENDS}), nó dùng \textbf{OpenFeign} — client REST khai báo bằng
interface, tự động tra cứu địa chỉ qua Eureka. Các Feign client quan trọng:
comment-service/story-service/feed-service/search-service gọi sang
post-service, user-service, group-service, fanpage-service.

Mỗi \path{@FeignClient} có kèm \textbf{Resilience4j circuit breaker} với
\path{fallbackFactory} riêng — phân biệt rõ hai loại lỗi: lỗi nghiệp vụ (ví dụ
404 không tìm thấy tài nguyên — vẫn trả lỗi đúng ngữ nghĩa cho client) và lỗi
hạ tầng (service đích sập/timeout — trả về \path{ServiceUnavailableException}
kèm mã 503 hoặc giá trị rỗng an toàn tuỳ ngữ cảnh, tránh một service sập kéo
sập cả chuỗi). Trạng thái circuit breaker có thể theo dõi qua
\path{actuator/circuitbreakers} của từng service.

\subsection{Giao tiếp bất đồng bộ — Apache Kafka}
\label{sec:kafka}

Các sự kiện không cần phản hồi ngay (fan-out thông báo, đồng bộ bộ đếm,
cascade xoá dữ liệu...) được phát qua Kafka, giúp service phát sự kiện không
phải chờ (và không phụ thuộc trực tiếp vào) mọi service tiêu thụ. Bảng dưới
liệt kê toàn bộ topic đã xác định trong mã nguồn (\path{KafkaTopics}):

\begin{longtable}{p{3.2cm}p{3.5cm}p{7.3cm}}
\threecolTableHeader{Bảng sự kiện Kafka}{Topic}{Producer}{Consumer(s) \& mục đích}
\path{USER_REGISTERED} & auth-service & user-service (\path{UserRegisteredListener}) — tự tạo hồ sơ (\path{UserProfile}) rỗng ngay sau khi tài khoản được tạo. \\
\path{POST_CREATED} & post-service & feed-service (\path{PostCreatedConsumer}) — fan-out bài viết mới vào Redis feed của bạn bè/người theo dõi. \\
\path{POST_TAGGED} & post-service & notification-service — tạo thông báo \path{TAG} cho người được gắn thẻ. \\
\path{COMMENT_CREATED} & comment-service & notification-service (thông báo cho chủ bài viết); reels-service (\path{ReelEventListener}, đồng bộ \path{commentCount} nếu đối tượng là reel). \\
\path{REACTION} & reaction-service & notification-service (thông báo); reels-service (đồng bộ \path{reactionCount}). \\
\path{GROUP} & group-service & notification-service — thông báo mời/duyệt thành viên nhóm. \\
\path{PAGE} & fanpage-service & (dự trữ cho mở rộng thông báo liên quan trang). \\
\path{FRIEND_REQUEST} & user-service (\path{FriendshipService}) & notification-service — thông báo lời mời kết bạn/chấp nhận. \\
\path{FOLLOW} & user-service (\path{FollowService}) & notification-service — thông báo có người theo dõi mới. \\
\path{MATCH} & dating-service & notification-service (thông báo ghép đôi thành công); chat-service (\path{MatchEventListener}) — tự tạo cuộc trò chuyện giữa hai người vừa match. \\
\path{MESSAGE} & chat-service & notification-service — thông báo tin nhắn mới (khi người nhận offline). \\
\path{REEL_CREATED} & reels-service & (dự trữ cho mở rộng, ví dụ fan-out feed reels). \\
\path{STORY_CREATED} & story-service & (dự trữ cho mở rộng). \\
\path{CONTENT_REMOVED} & moderation-service & comment-service, group-service, fanpage-service, reels-service, story-service — xoá cascade nội dung/dữ liệu liên quan khi quản trị viên duyệt gỡ vi phạm. \\
\end{longtable}

\section{Lưu trữ dữ liệu — Database per Service}

Mỗi service nghiệp vụ sở hữu cơ sở dữ liệu riêng, không service nào truy cập
trực tiếp database của service khác (đúng nguyên lý \emph{database-per-service}
của microservices) — mọi truy cập chéo đều phải qua API (Feign) hoặc sự kiện
(Kafka).

\begin{longtable}{p{3.5cm}p{3cm}p{8cm}}
\threecolTableHeader{Bảng cơ sở dữ liệu theo service}{Service}{Loại DB}{Tên DB / namespace}
auth-service & PostgreSQL & \path{auth_db} \\
user-service & PostgreSQL & \path{user_db} \\
media-service & PostgreSQL + MinIO & \path{media_db} (metadata) + bucket đối tượng nhị phân \\
post-service & PostgreSQL & \path{post_db} \\
comment-service & PostgreSQL & \path{comment_db} \\
reaction-service & PostgreSQL & \path{reaction_db} \\
group-service & PostgreSQL & \path{group_db} \\
fanpage-service & PostgreSQL & \path{fanpage_db} \\
dating-service & PostgreSQL & \path{dating_db} \\
moderation-service & PostgreSQL & \path{moderation_db} \\
story-service & MongoDB & \path{story_db} \\
reels-service & MongoDB & \path{reels_db} \\
chat-service & MongoDB + Redis & \path{chat_db} (hội thoại/tin nhắn) + Redis (trạng thái online) \\
notification-service & MongoDB & \path{notification_db} (TTL 8 tuần) \\
feed-service & Redis & Sorted set theo từng user (không có schema quan hệ) \\
search-service & Không có & Tổng hợp qua Feign, không lưu trữ dữ liệu riêng \\
\end{longtable}

10 database PostgreSQL được khởi tạo tự động trong cùng một instance
PostgreSQL (container \path{sma-postgres}) thông qua script
\path{docker/postgres/init-multi-db.sh} chạy lúc container khởi động lần đầu
(biến môi trường \path{POSTGRES_MULTIPLE_DATABASES}).

% ============================================================
\chapter{Xác thực và bảo mật}
\label{ch:security}
% ============================================================

\section{Luồng đăng ký / đăng nhập}

\path{auth-service} (cổng 8081, database \path{auth_db}) là nơi duy nhất phát
hành JWT trong toàn hệ thống:

\begin{enumerate}[nosep]
    \item \path{POST /api/auth/register} — tạo \path{Account} mới (mật khẩu
    băm bằng BCrypt), mặc định role \path{USER}; nếu email nằm trong danh sách
    \path{ADMIN_EMAILS} (biến môi trường, phân tách bằng dấu phẩy) thì tự động
    gán thêm role \path{ADMIN}. Sau khi tạo thành công, phát sự kiện Kafka
    \path{USER_REGISTERED} để user-service tự tạo hồ sơ tương ứng.
    \item \path{POST /api/auth/login} — xác thực email/mật khẩu, phát hành cặp
    \emph{access token} (JWT, hạn 15 phút — \path{JWT_ACCESS_EXPIRATION_MS}
    mặc định \path{900000}) và \emph{refresh token} (JWT, hạn 7 ngày —
    \path{JWT_REFRESH_EXPIRATION_MS} mặc định \path{604800000}, lưu vào bảng
    \path{RefreshToken} để có thể thu hồi).
    \item \path{POST /api/auth/refresh} — đổi refresh token còn hạn lấy access
    token mới mà không cần đăng nhập lại.
    \item \path{POST /api/auth/logout} — thu hồi refresh token hiện tại.
    \item \path{GET /api/auth/me} — trả thông tin tài khoản đang đăng nhập
    (dựa trên header \path{X-User-Id} gateway đã gắn).
\end{enumerate}

\section{Xác thực tập trung tại API Gateway}

Toàn bộ việc kiểm tra JWT được tập trung ở một \path{GlobalFilter} duy nhất
của gateway — \path{JwtAuthenticationFilter} — thay vì lặp lại logic này ở
từng service:

\begin{enumerate}[nosep]
    \item Nếu đường dẫn request khớp một trong các mẫu công khai
    (\path{PUBLIC_PATTERNS}: \path{/api/auth/register}, \path{/api/auth/login},
    \path{/api/auth/refresh}, \path{/actuator/**}, \path{/eureka/**},
    \path{/swagger-ui.html}, \path{/swagger-ui/**}, \path{/webjars/**},
    \path{/v3/api-docs/**}, \path{/docs/**}) — cho qua thẳng, không cần token.
    \item Ngược lại, trích token từ header \path{Authorization: Bearer <token>};
    nếu không có, thử lấy từ query param \path{?token=} — dành riêng cho
    WebSocket client (trình duyệt không phải lúc nào cũng gắn được custom
    header vào request nâng cấp giao thức WebSocket).
    \item Token không hợp lệ/hết hạn → trả về ngay \path{401 Unauthorized}
    (JSON \path{ApiResponse.error(...)}), không forward xuống service phía sau.
    \item Token hợp lệ → giải mã lấy \path{userId} và danh sách \path{roles},
    gắn vào request dưới dạng header \path{X-User-Id} và \path{X-User-Roles}
    (phân tách bằng dấu phẩy) rồi mới forward xuống business service. Từ đây,
    \textbf{business service không bao giờ tự parse JWT} — chỉ đọc hai header
    này để biết danh tính người gọi.
\end{enumerate}

\section{Phân quyền theo vai trò (Role-based access)}

Hệ thống có hai vai trò: \path{USER} (mặc định) và \path{ADMIN}. Vai trò
\path{ADMIN} chỉ được cấp tự động cho các email nằm trong allowlist
\path{ADMIN_EMAILS} lúc đăng ký (không có API tự nâng quyền). Vai trò này được
dùng để gate các API quản trị nội dung tại \path{moderation-service}
(\path{GET /api/moderation/reports}, \path{PUT /api/moderation/reports/\{id\}/resolve})
và trang quản trị trên frontend.

\section{Giới hạn tần suất (Rate limiting)}

Ngoài xác thực, gateway còn áp \path{RequestRateLimiter} (thuật toán token
bucket, trạng thái lưu trên Redis) cho từng route — xem bảng định tuyến ở
mục~2.2. Khoá giới hạn (\path{key-resolver}) là \path{userOrIpKeyResolver}:
ưu tiên định danh theo \path{userId} (đã xác thực) để tránh một người dùng
dùng nhiều IP né giới hạn, và rơi về IP nguồn khi request chưa xác thực (ví dụ
chính các endpoint đăng ký/đăng nhập). Route \path{/api/auth/**} bị giới hạn
chặt hơn hẳn các route khác (1 req/s, burst 5) để giảm rủi ro brute-force mật
khẩu; hai route WebSocket cũng bị giới hạn số lần \emph{thiết lập kết nối}
(5 req/s, burst 10) để tránh mở tràn lan connection.

% ============================================================
\chapter{Đặc tả chi tiết các microservice nghiệp vụ}
\label{ch:services}
% ============================================================

Chương này đặc tả từng business service: mục đích, mô hình dữ liệu (thực thể
và các trường chính) và danh sách API tiêu biểu (đường dẫn luôn được viết theo
dạng đã đi qua API Gateway, ví dụ \path{/api/posts/...}, vì đây là cách duy
nhất client bên ngoài được phép gọi vào hệ thống).

\section{auth-service (cổng 8081)}

\textbf{Mục đích:} phát hành và quản lý vòng đời JWT (đăng ký, đăng nhập, làm
mới, đăng xuất) — luồng chi tiết đã trình bày ở Chương~\ref{ch:security}.

\subsubsection*{Cơ sở dữ liệu (\path{auth_db})}
\begin{longtable}{p{3.2cm}p{3cm}p{8.3cm}}
\entityTableHeader{Account}
\multicolumn{3}{l}{\textbf{Account} (bảng \path{accounts})} \\
\path{id} & \path{String (UUID)} & Khoá chính. \\
\path{email} & \path{String} & Duy nhất, dùng để đăng nhập. \\
\path{passwordHash} & \path{String} & Mật khẩu đã băm (BCrypt). \\
\path{phone} & \path{String} & Không bắt buộc. \\
\path{status} & \path{Enum} & \path{ACTIVE} mặc định. \\
\path{roles} & \path{List<String>} & Bảng phụ \path{account_roles}; mặc định \path{["USER"]}. \\
\path{fullName} & \path{String} & \\
\path{createdAt} & \path{Instant} & \\
\midrule
\multicolumn{3}{l}{\textbf{RefreshToken} (bảng \path{refresh_tokens})} \\
\path{id} & \path{String (UUID)} & Khoá chính. \\
\path{accountId} & \path{String} & Tham chiếu \path{Account.id}. \\
\path{token} & \path{String} & Chuỗi JWT refresh, duy nhất. \\
\path{expiresAt} & \path{Instant} & \\
\path{revoked} & \path{boolean} & Đặt \path{true} khi logout hoặc bị thu hồi. \\
\end{longtable}

\subsubsection*{API chính}
\begin{longtable}{p{2cm}p{5cm}p{7.5cm}}
\apiTableHeader{auth-service}
\path{POST} & \path{/api/auth/register} & Đăng ký tài khoản mới. Công khai. \\
\path{POST} & \path{/api/auth/login} & Đăng nhập, trả access + refresh token. Công khai. \\
\path{POST} & \path{/api/auth/refresh} & Đổi refresh token lấy access token mới. Công khai. \\
\path{POST} & \path{/api/auth/logout} & Thu hồi refresh token hiện tại. \\
\path{GET} & \path{/api/auth/me} & Thông tin tài khoản đang đăng nhập. \\
\end{longtable}

\section{user-service (cổng 8082)}

\textbf{Mục đích:} hồ sơ người dùng, quan hệ xã hội (kết bạn, theo dõi, chặn)
và gợi ý kết bạn.

\subsubsection*{Cơ sở dữ liệu (\path{user_db})}
\begin{longtable}{p{3.2cm}p{3cm}p{8.3cm}}
\entityTableHeader{user-service}
\multicolumn{3}{l}{\textbf{UserProfile} (bảng \path{user_profiles})} \\
\path{id} & \path{String} & Bằng \path{Account.id} bên auth-service (tạo tự động qua sự kiện \path{USER_REGISTERED}). \\
\path{fullName}, \path{avatarUrl}, \path{coverUrl}, \path{bio} & & \\
\path{dob} & \path{LocalDate} & Ngày sinh. \\
\path{gender} & \path{Enum} & \\
\path{location} & \path{String} & \\
\path{workplace} & \path{String} & Dòng tự do kiểu "Làm việc tại X / Học tại Y" (không mô hình hoá đa lịch sử công việc như Facebook thật). \\
\path{readReceiptsEnabled} & \path{boolean} & Bật/tắt "đã xem" hai chiều cho chat — tắt cũng ẩn luôn trạng thái đã xem của người khác với mình. \\
\midrule
\multicolumn{3}{l}{\textbf{Friendship} (bảng \path{friendships}, unique theo cặp requester/addressee)} \\
\path{requesterId}, \path{addresseeId} & \path{String} & \\
\path{status} & \path{Enum} & \path{PENDING}/\path{ACCEPTED}/\path{DECLINED}. \\
\path{createdAt}, \path{respondedAt} & \path{Instant} & \\
\midrule
\multicolumn{3}{l}{\textbf{Follow} (bảng \path{follows}, unique theo cặp follower/followee)} \\
\path{followerId}, \path{followeeId} & \path{String} & \\
\path{createdAt} & \path{Instant} & \\
\midrule
\multicolumn{3}{l}{\textbf{Block} (bảng \path{blocks}, unique theo cặp blocker/blocked)} \\
\path{blockerId}, \path{blockedId} & \path{String} & \\
\path{createdAt} & \path{Instant} & \\
\midrule
\multicolumn{3}{l}{\textbf{FriendSuggestionDismissal}} \\
\multicolumn{3}{p{14.5cm}}{Ghi nhận gợi ý kết bạn đã bị người dùng bỏ qua, để không đề xuất lại.} \\
\end{longtable}

\subsubsection*{API chính}
\begin{longtable}{p{2cm}p{5.5cm}p{7cm}}
\apiTableHeader{user-service}
\path{GET} & \path{/api/users/\{id\}} & Xem hồ sơ người dùng khác. \\
\path{GET} & \path{/api/users/search} & Tìm người dùng theo tên. \\
\path{GET} & \path{/api/users/me} & Hồ sơ của chính mình. \\
\path{PUT} & \path{/api/users/me} & Cập nhật hồ sơ. \\
\path{PUT} & \path{/api/users/me/avatar} & Cập nhật ảnh đại diện. \\
\path{PUT} & \path{/api/users/me/cover} & Cập nhật ảnh bìa. \\
\path{POST} & \path{/api/users/\{targetId\}/follow} & Theo dõi người dùng. \\
\path{DELETE} & \path{/api/users/\{targetId\}/follow} & Bỏ theo dõi. \\
\path{GET} & \path{/api/users/\{targetId\}/follow-status} & Kiểm tra trạng thái theo dõi. \\
\path{POST} & \path{/api/users/\{targetId\}/friend-request} & Gửi lời mời kết bạn. \\
\path{PUT} & \path{/api/users/friend-requests/\{id\}/accept} & Chấp nhận lời mời. \\
\path{PUT} & \path{/api/users/friend-requests/\{id\}/decline} & Từ chối lời mời. \\
\path{DELETE} & \path{/api/users/friends/\{friendId\}} & Huỷ kết bạn. \\
\path{GET} & \path{/api/users/me/friends} & Danh sách bạn bè. \\
\path{GET} & \path{/api/users/me/friend-requests} & Lời mời kết bạn đang chờ. \\
\path{GET} & \path{/api/users/\{targetId\}/friendship-status} & Trạng thái quan hệ bạn bè với một người. \\
\path{GET} & \path{/api/users/\{id\}/friend-ids} & Danh sách ID bạn bè (dùng nội bộ, ví dụ post-service kiểm tra quyền \path{FRIENDS}). \\
\path{POST}/\path{DELETE} & \path{/api/users/\{targetId\}/block} & Chặn / bỏ chặn người dùng. \\
\path{GET} & \path{/api/users/me/blocked} & Danh sách đã chặn. \\
\path{DELETE} & \path{/api/users/friend-suggestions/\{targetId\}} & Bỏ qua một gợi ý kết bạn. \\
\end{longtable}

\section{media-service (cổng 8083)}

\textbf{Mục đích:} nhận upload và cấp phát URL cho file media (ảnh/video),
lưu nhị phân trên MinIO (tương thích S3), lưu metadata trên PostgreSQL.

\subsubsection*{Cơ sở dữ liệu (\path{media_db})}
\begin{longtable}{p{3.2cm}p{3cm}p{8.3cm}}
\entityTableHeader{MediaFile}
\path{id} & \path{String (UUID)} & \\
\path{ownerId} & \path{String} & Người upload. \\
\path{objectKey} & \path{String} & Khoá đối tượng trong bucket MinIO. \\
\path{url} & \path{String} & URL truy cập công khai/đã ký. \\
\path{contentType} & \path{String} & MIME type. \\
\path{purpose} & \path{Enum (MediaPurpose)} & Ví dụ: \path{AVATAR}, \path{COVER}, \path{POST}, \path{STORY}, \path{REEL}, ... \\
\path{sizeBytes} & \path{long} & \\
\path{createdAt} & \path{Instant} & \\
\end{longtable}

\subsubsection*{API chính}
\begin{longtable}{p{2cm}p{5cm}p{7.5cm}}
\apiTableHeader{media-service}
\path{POST} & \path{/api/media} & Upload file (multipart), trả về metadata + URL. \\
\path{GET} & \path{/api/media/\{id\}} & Lấy metadata một file. \\
\path{DELETE} & \path{/api/media/\{id\}} & Xoá file (cả object trên MinIO lẫn metadata). \\
\end{longtable}

\section{post-service (cổng 8084)}

\textbf{Mục đích:} CRUD bài viết, chia sẻ/repost, gắn thẻ người dùng, quyền
riêng tư (kể cả tuỳ biến theo danh sách người xem cụ thể), ghim bài, đăng bài
trong nhóm/trang.

\subsubsection*{Cơ sở dữ liệu (\path{post_db})}
\begin{longtable}{p{3.2cm}p{3cm}p{8.3cm}}
\entityTableHeader{Post}
\path{id} & \path{String (UUID)} & \\
\path{authorId} & \path{String} & \\
\path{content} & \path{Text} & \\
\path{mediaUrls} & \path{List<String>} & Bảng phụ \path{post_media_urls}. \\
\path{privacy} & \path{Enum} & \path{PUBLIC}/\path{FRIENDS}/\path{CUSTOM}/... \\
\path{customAudienceUserIds} & \path{List<String>} & Chỉ dùng khi \path{privacy = CUSTOM}; bảng phụ \path{post_custom_audience}. \\
\path{taggedUserIds} & \path{List<String>} & Bảng phụ \path{post_tagged_users}; phát sự kiện \path{POST_TAGGED} cho từng người được gắn thẻ. \\
\path{groupId}, \path{pageId} & \path{String} & Khác null nếu bài đăng trong nhóm/trang. \\
\path{sharedPostId} & \path{String} & Khác null nếu đây là bài chia sẻ lại một bài viết khác (lấy riêng qua \path{GET /api/posts/\{sharedPostId\}}). \\
\path{sharedReelId} & \path{String} & Khác null nếu chia sẻ lại một reel (loại trừ lẫn nhau với \path{sharedPostId}). \\
\path{commentCount}, \path{reactionCount}, \path{shareCount} & \path{int} & Bộ đếm, đồng bộ qua Kafka. \\
\path{pinned} & \path{boolean} & Bài viết được ghim lên đầu trang cá nhân. \\
\path{createdAt}, \path{updatedAt} & \path{Instant} & \\
\end{longtable}

\subsubsection*{API chính}
\begin{longtable}{p{2cm}p{5.5cm}p{7cm}}
\apiTableHeader{post-service}
\path{POST} & \path{/api/posts} & Tạo bài viết mới; phát sự kiện \path{POST_CREATED} (và \path{POST_TAGGED} nếu có gắn thẻ). \\
\path{GET} & \path{/api/posts/\{id\}} & Xem chi tiết một bài viết (kiểm tra quyền riêng tư). \\
\path{GET} & \path{/api/posts/search} & Tìm bài viết theo nội dung. \\
\path{PUT} & \path{/api/posts/\{id\}} & Sửa bài viết (chỉ tác giả). \\
\path{DELETE} & \path{/api/posts/\{id\}} & Xoá bài viết. \\
\path{GET} & \path{/api/posts/author/\{authorId\}} & Danh sách bài viết của một người (enforce quyền riêng tư theo người xem hiện tại, có gọi Feign sang user-service kiểm tra bạn bè nếu \path{privacy = FRIENDS}). \\
\path{PUT}/\path{DELETE} & \path{/api/posts/\{id\}/pin} & Ghim / bỏ ghim bài viết. \\
\path{POST} & \path{/api/posts/\{id\}/share} & Chia sẻ lại một bài viết (tăng \path{shareCount} bài gốc). \\
\path{POST} & \path{/api/posts/share-reel/\{reelId\}} & Chia sẻ lại một reel dưới dạng bài viết. \\
\path{GET} & \path{/api/posts/group/\{groupId\}} & Bài viết trong một nhóm. \\
\path{GET} & \path{/api/posts/page/\{pageId\}} & Bài viết trong một trang (fanpage). \\
\path{GET} & \path{/api/posts/batch} & Lấy nhiều bài viết theo danh sách ID (dùng nội bộ, ví dụ feed-service/search-service). \\
\end{longtable}

\section{comment-service (cổng 8085)}

\textbf{Mục đích:} bình luận (kèm trả lời lồng nhau) trên bài viết hoặc reel.

\subsubsection*{Cơ sở dữ liệu (\path{comment_db})}
\begin{longtable}{p{3.2cm}p{3cm}p{8.3cm}}
\entityTableHeader{Comment}
\path{id} & \path{String (UUID)} & \\
\path{targetType} & \path{Enum (TargetType)} & \path{POST} hoặc \path{REEL} — cùng cấu trúc với reaction-service để dễ mở rộng thêm loại nội dung khác sau này. \\
\path{targetId} & \path{String} & ID bài viết/reel được bình luận. \\
\path{authorId} & \path{String} & \\
\path{targetOwnerId} & \path{String} & Chủ sở hữu nội dung (để phát thông báo mà không cần gọi ngược sang post/reels-service). \\
\path{parentCommentId} & \path{String} & Khác null nếu là trả lời một bình luận khác. \\
\path{content} & \path{Text} & \\
\path{deleted} & \path{boolean} & Xoá mềm. \\
\path{createdAt}, \path{updatedAt} & \path{Instant} & \\
\end{longtable}

\subsubsection*{API chính}
\begin{longtable}{p{2cm}p{5.5cm}p{7cm}}
\apiTableHeader{comment-service}
\path{POST} & \path{/api/comments} & Tạo bình luận mới; phát sự kiện \path{COMMENT_CREATED}. \\
\path{GET} & \path{/api/comments/\{id\}/replies} & Danh sách trả lời của một bình luận. \\
\path{PUT} & \path{/api/comments/\{id\}} & Sửa nội dung bình luận. \\
\path{DELETE} & \path{/api/comments/\{id\}} & Xoá (mềm) bình luận. \\
\end{longtable}

\section{reaction-service (cổng 8086)}

\textbf{Mục đích:} quản lý cảm xúc (reaction) và mục đã lưu (saved item) trên
bài viết/bình luận/reel — hai chức năng gộp chung một service vì cùng mô hình
truy cập dữ liệu.

\subsubsection*{Cơ sở dữ liệu (\path{reaction_db})}
\begin{longtable}{p{3.2cm}p{3cm}p{8.3cm}}
\entityTableHeader{reaction-service}
\multicolumn{3}{l}{\textbf{Reaction} (bảng \path{reactions}, unique theo (targetType, targetId, userId))} \\
\path{targetType} & \path{Enum} & Loại đối tượng được react. \\
\path{targetId} & \path{String} & \\
\path{targetOwnerId} & \path{String} & Do client cung cấp trực tiếp (đơn giản hoá có chủ đích — xem Chương~\ref{ch:limitations}). \\
\path{userId} & \path{String} & Người react. \\
\path{type} & \path{Enum (ReactionType)} & LIKE/LOVE/HAHA/... \\
\path{createdAt} & \path{Instant} & \\
\midrule
\multicolumn{3}{l}{\textbf{SavedItem} (bảng \path{saved_items}, unique theo (targetType, targetId, userId))} \\
\multicolumn{3}{p{14.5cm}}{Cấu trúc giống hệt Reaction (không có trường \path{type}) — đại diện cho việc một người dùng "lưu" một nội dung.} \\
\end{longtable}

\subsubsection*{API chính}
\begin{longtable}{p{2cm}p{5.5cm}p{7cm}}
\apiTableHeader{reaction-service}
\path{POST} & \path{/api/reactions} & Thả/đổi cảm xúc lên một đối tượng; phát sự kiện Kafka \path{REACTION}. \\
\path{GET} & \path{/api/reactions/summary} & Tổng hợp số lượng theo từng loại cảm xúc cho một đối tượng. \\
\path{GET} & \path{/api/reactions/me} & Cảm xúc hiện tại của người dùng trên một đối tượng. \\
\path{POST}/\path{DELETE} & \path{/api/saved} & Lưu / bỏ lưu một đối tượng. \\
\path{GET} & \path{/api/saved/me} & Kiểm tra một đối tượng có đang được lưu không. \\
\path{GET} & \path{/api/saved/mine} & Danh sách toàn bộ mục đã lưu của người dùng. \\
\end{longtable}

\section{story-service (cổng 8087)}

\textbf{Mục đích:} story dạng ảnh/video tự động biến mất sau một khoảng thời
gian (mặc định theo Facebook: 24 giờ).

\subsubsection*{Cơ sở dữ liệu (MongoDB, \path{story_db})}
\begin{longtable}{p{3.2cm}p{3cm}p{8.3cm}}
\entityTableHeader{Story}
\path{id} & \path{String} & \\
\path{authorId} & \path{String} & \\
\path{mediaUrl}, \path{mediaType} & & Ảnh hoặc video. \\
\path{caption} & \path{String} & \\
\path{viewerIds} & \path{Set<String>} & Người đã xem. \\
\path{textOverlays} & \path{List<TextOverlay>} & Văn bản chèn lên story (nội dung, vị trí, kiểu chữ). \\
\path{createdAt}, \path{expiresAt} & \path{Instant} & \path{expiresAt} có \emph{TTL index} (\path{expireAfterSeconds = 0}) — MongoDB tự xoá document đúng thời điểm này mà không cần job dọn dẹp riêng. \\
\end{longtable}

\subsubsection*{API chính}
\begin{longtable}{p{2cm}p{5.5cm}p{7cm}}
\apiTableHeader{story-service}
\path{POST} & \path{/api/stories} & Đăng story mới; phát sự kiện \path{STORY_CREATED}. \\
\path{GET} & \path{/api/stories/feed} & Story còn hiệu lực của bạn bè/người theo dõi. \\
\path{GET} & \path{/api/stories/author/\{authorId\}} & Story còn hiệu lực của một người cụ thể. \\
\path{POST} & \path{/api/stories/\{id\}/view} & Đánh dấu đã xem. \\
\path{DELETE} & \path{/api/stories/\{id\}} & Xoá story trước hạn. \\
\end{longtable}

\section{reels-service (cổng 8088)}

\textbf{Mục đích:} video ngắn dạng Reels, gồm bộ đếm lượt xem/bình luận/cảm
xúc/chia sẻ được đồng bộ từ các service khác qua Kafka.

\subsubsection*{Cơ sở dữ liệu (MongoDB, \path{reels_db})}
\begin{longtable}{p{3.2cm}p{3cm}p{8.3cm}}
\entityTableHeader{Reel}
\path{id} & \path{String} & \\
\path{authorId} & \path{String} & \\
\path{videoUrl}, \path{thumbnailUrl} & \path{String} & \\
\path{caption} & \path{String} & \\
\path{viewCount} & \path{long} & \\
\path{commentCount}, \path{reactionCount} & \path{int} & Đồng bộ từ sự kiện \path{COMMENT_CREATED}/\path{REACTION} qua \path{ReelEventListener}. \\
\path{shareCount} & \path{int} & \\
\path{createdAt} & \path{Instant} & \\
\end{longtable}

\subsubsection*{API chính}
\begin{longtable}{p{2cm}p{5.5cm}p{7cm}}
\apiTableHeader{reels-service}
\path{POST} & \path{/api/reels} & Đăng reel mới; phát sự kiện \path{REEL_CREATED}. \\
\path{GET} & \path{/api/reels/\{id\}} & Xem chi tiết. \\
\path{POST} & \path{/api/reels/\{id\}/view} & Tăng lượt xem. \\
\path{POST} & \path{/api/reels/\{id\}/share-count} & Tăng bộ đếm chia sẻ (gọi từ post-service khi share-reel). \\
\path{GET} & \path{/api/reels/feed} & Danh sách reel dạng cuộn (feed). \\
\path{GET} & \path{/api/reels/author/\{authorId\}} & Reel của một người. \\
\path{DELETE} & \path{/api/reels/\{id\}} & Xoá reel. \\
\end{longtable}

\section{group-service (cổng 8089)}

\textbf{Mục đích:} nhóm (group) công khai/riêng tư, quản lý thành viên và vai
trò trong nhóm.

\subsubsection*{Cơ sở dữ liệu (\path{group_db})}
\begin{longtable}{p{3.2cm}p{3cm}p{8.3cm}}
\entityTableHeader{group-service}
\multicolumn{3}{l}{\textbf{Group} (entity JPA tên \path{SocialGroup}, bảng \path{groups})} \\
\path{id} & \path{String (UUID)} & \\
\path{name}, \path{description} & & \\
\path{avatarUrl}, \path{coverUrl} & & \\
\path{privacy} & \path{Enum} & \path{PUBLIC} / \path{PRIVATE}. \\
\path{ownerId} & \path{String} & \\
\path{memberCount} & \path{int} & Mặc định 1 (chủ nhóm). \\
\path{createdAt} & \path{Instant} & \\
\midrule
\multicolumn{3}{l}{\textbf{GroupMember} (bảng \path{group_members}, unique theo (groupId, userId))} \\
\path{groupId}, \path{userId} & \path{String} & \\
\path{role} & \path{Enum (MemberRole)} & Ví dụ \path{OWNER}/\path{ADMIN}/\path{MEMBER}. \\
\path{status} & \path{Enum (MemberStatus)} & Ví dụ \path{PENDING}/\path{APPROVED}. \\
\path{joinedAt} & \path{Instant} & \\
\end{longtable}

\subsubsection*{API chính}
\begin{longtable}{p{2.4cm}p{5.5cm}p{6.6cm}}
\apiTableHeader{group-service}
\path{POST} & \path{/api/groups} & Tạo nhóm mới. \\
\path{GET} & \path{/api/groups} & Danh sách nhóm — bao gồm cả nhóm \path{PRIVATE} mà người gọi đã là thành viên \path{APPROVED}, không chỉ nhóm \path{PUBLIC}. \\
\path{GET} & \path{/api/groups/\{id\}} & Chi tiết một nhóm. \\
\path{POST} & \path{/api/groups/\{id\}/join} & Xin gia nhập nhóm (vào \path{PENDING} nếu nhóm riêng tư). \\
\path{PUT} & \path{/api/groups/\{id\}/members/\{userId\}/approve} & Duyệt thành viên (quản trị nhóm). \\
\path{DELETE} & \path{/api/groups/\{id\}/leave} & Rời nhóm. \\
\path{DELETE} & \path{/api/groups/\{id\}/members/\{userId\}} & Kick thành viên. \\
\path{PUT} & \path{/api/groups/\{id\}/members/\{userId\}/role} & Đổi vai trò thành viên. \\
\path{GET} & \path{/api/groups/\{id\}/members} & Danh sách thành viên. \\
\path{GET} & \path{/api/groups/me} & Danh sách nhóm mình tham gia. \\
\end{longtable}

\textbf{Ghi chú:} \path{GroupModerationListener} lắng nghe sự kiện
\path{CONTENT_REMOVED} từ moderation-service để xoá cascade nhóm cùng toàn bộ
bảng thành viên liên quan khi quản trị viên duyệt gỡ vi phạm.

\section{fanpage-service (cổng 8090)}

\textbf{Mục đích:} trang (fanpage) — luôn công khai theo thiết kế (không có
khái niệm trang riêng tư như nhóm), có nhiều quản trị viên và người theo dõi.

\subsubsection*{Cơ sở dữ liệu (\path{fanpage_db})}
\begin{longtable}{p{3.2cm}p{3cm}p{8.3cm}}
\entityTableHeader{fanpage-service}
\multicolumn{3}{l}{\textbf{Fanpage} (bảng \path{fanpages})} \\
\path{id} & \path{String (UUID)} & \\
\path{name}, \path{category}, \path{description} & & \\
\path{avatarUrl}, \path{coverUrl} & & \\
\path{ownerId} & \path{String} & \\
\path{followerCount} & \path{int} & \\
\path{createdAt} & \path{Instant} & \\
\midrule
\multicolumn{3}{l}{\textbf{PageAdmin} (bảng \path{page_admins}, unique theo (pageId, userId))} \\
\path{pageId}, \path{userId} & \path{String} & \\
\path{role} & \path{Enum (AdminRole)} & \\
\path{addedAt} & \path{Instant} & \\
\midrule
\multicolumn{3}{l}{\textbf{PageFollower} (bảng \path{page_followers}, unique theo (pageId, userId))} \\
\path{pageId}, \path{userId} & \path{String} & \\
\path{createdAt} & \path{Instant} & \\
\end{longtable}

\subsubsection*{API chính}
\begin{longtable}{p{2cm}p{5.5cm}p{7cm}}
\apiTableHeader{fanpage-service}
\path{POST} & \path{/api/pages} & Tạo trang mới. \\
\path{GET} & \path{/api/pages/\{id\}} & Chi tiết một trang. \\
\path{POST}/\path{DELETE} & \path{/api/pages/\{id\}/follow} & Theo dõi / bỏ theo dõi trang. \\
\path{POST} & \path{/api/pages/\{id\}/admins} & Thêm quản trị viên trang. \\
\path{DELETE} & \path{/api/pages/\{id\}/admins/\{userId\}} & Gỡ quản trị viên. \\
\path{GET} & \path{/api/pages/\{id\}/followers} & Danh sách người theo dõi. \\
\path{GET} & \path{/api/pages/me/managed} & Trang mình đang quản trị. \\
\path{GET} & \path{/api/pages/me/followed} & Trang mình đang theo dõi. \\
\end{longtable}

\section{dating-service (cổng 8091)}

\textbf{Mục đích:} tính năng hẹn hò kiểu Tinder — hồ sơ hẹn hò riêng (khác hồ
sơ mạng xã hội), vuốt (swipe) và ghép đôi (match) theo thuật toán chấm điểm
(chi tiết thuật toán ở mục~\ref{sec:matching}).

\subsubsection*{Cơ sở dữ liệu (\path{dating_db})}
\begin{longtable}{p{3.2cm}p{3cm}p{8.3cm}}
\entityTableHeader{dating-service}
\multicolumn{3}{l}{\textbf{DatingProfile} (bảng \path{dating_profiles})} \\
\path{id} & \path{String} & Bằng \path{userId}. \\
\path{gender} & \path{Enum} & Giới tính của chính chủ hồ sơ. \\
\path{birthDate} & \path{LocalDate} & \\
\path{bio} & \path{Text} & \\
\path{interests} & \path{List<String>} & Bảng phụ \path{dating_profile_interests}. \\
\path{minAgePreference}, \path{maxAgePreference} & \path{int} & \\
\path{genderPreference} & \path{Enum} & Giới tính \emph{đang tìm} — khác với trường \path{gender} ở trên. \\
\path{photos} & \path{List<String>} & Bảng phụ \path{dating_profile_photos}. \\
\path{active} & \path{boolean} & Tắt để tạm ẩn khỏi vòng quay ghép đôi. \\
\path{createdAt} & \path{Instant} & \\
\midrule
\multicolumn{3}{l}{\textbf{Swipe} (bảng \path{swipes}, unique theo (swiperId, targetId))} \\
\path{swiperId}, \path{targetId} & \path{String} & \\
\path{action} & \path{Enum (SwipeAction)} & \path{LIKE} / \path{PASS}. \\
\path{createdAt} & \path{Instant} & \\
\midrule
\multicolumn{3}{l}{\textbf{Match} (bảng \path{matches}, unique theo cặp (user1Id, user2Id))} \\
\multicolumn{3}{p{14.5cm}}{Được tạo khi cả hai phía đều \path{LIKE} nhau (mutual like); kèm phát sự kiện Kafka \path{MATCH}.} \\
\end{longtable}

\subsubsection*{API chính}
\begin{longtable}{p{2cm}p{5.5cm}p{7cm}}
\apiTableHeader{dating-service}
\path{PUT} & \path{/api/dating/profile} & Tạo/cập nhật hồ sơ hẹn hò. \\
\path{GET} & \path{/api/dating/profile/me} & Xem hồ sơ hẹn hò của mình. \\
\path{GET} & \path{/api/dating/candidates} & Danh sách ứng viên phù hợp, kèm điểm số (0--100). \\
\path{POST} & \path{/api/dating/swipe} & Vuốt LIKE/PASS một ứng viên; nếu mutual LIKE thì tạo \path{Match} + phát sự kiện. \\
\path{GET} & \path{/api/dating/matches} & Danh sách đã match. \\
\path{DELETE} & \path{/api/dating/matches/\{id\}} & Huỷ một match (unmatch). \\
\end{longtable}

\section{chat-service (cổng 8092)}

\textbf{Mục đích:} nhắn tin thời gian thực (1--1 và có thể mở rộng nhóm),
trạng thái online (presence), tự động tạo cuộc trò chuyện khi có match hẹn hò
mới.

\subsubsection*{Cơ sở dữ liệu (MongoDB \path{chat_db} + Redis presence)}
\begin{longtable}{p{3.2cm}p{3cm}p{8.3cm}}
\entityTableHeader{chat-service}
\multicolumn{3}{l}{\textbf{Conversation} (collection \path{conversations})} \\
\path{id} & \path{String} & \\
\path{type} & \path{Enum (ConversationType)} & Ví dụ 1--1, nhóm. \\
\path{participantIds} & \path{List<String>} & \\
\path{lastMessagePreview}, \path{lastMessageAt} & & Hiển thị nhanh trong danh sách hội thoại. \\
\path{createdAt} & \path{Instant} & \\
\midrule
\multicolumn{3}{l}{\textbf{Message} (collection \path{messages})} \\
\path{id}, \path{conversationId}, \path{senderId} & \path{String} & \\
\path{content}, \path{mediaUrl} & & \\
\path{storyReplyId}, \path{storyReplyPreviewUrl} & \path{String} & Khác null nếu tin nhắn là trả lời một story. \\
\path{sentAt} & \path{Instant} & \\
\path{readBy} & \path{Set<String>} & Người đã đọc (đọc-đã-xem hai chiều, xem \path{readReceiptsEnabled} ở user-service). \\
\path{deleted} & \path{boolean} & Thu hồi tin nhắn (xoá mềm). \\
\end{longtable}

\subsubsection*{API chính (REST)}
\begin{longtable}{p{2cm}p{5.5cm}p{7cm}}
\apiTableHeader{chat-service}
\path{POST} & \path{/api/chat/conversations} & Tạo (hoặc lấy đã có) hội thoại với một người. \\
\path{GET} & \path{/api/chat/conversations} & Danh sách hội thoại của mình. \\
\path{GET} & \path{/api/chat/conversations/\{id\}/messages} & Lịch sử tin nhắn (phân trang). \\
\path{POST} & \path{/api/chat/conversations/\{id\}/read} & Đánh dấu đã đọc. \\
\path{DELETE} & \path{/api/chat/messages/\{id\}} & Thu hồi tin nhắn. \\
\path{GET} & \path{/api/chat/presence/\{userId\}} & Trạng thái online của một người. \\
\path{POST} & \path{/api/chat/presence/batch} & Trạng thái online của nhiều người cùng lúc. \\
\end{longtable}

\subsubsection*{Kênh thời gian thực (WebSocket/STOMP)}
Gateway định tuyến riêng \path{/ws/**} (\path{uri: lb:ws://CHAT-SERVICE}) sang
chat-service. Client kết nối \path{ws://.../ws?token=<JWT>} (SockJS) hoặc
\path{ws://.../ws/websocket?token=<JWT>} (raw WebSocket, bỏ qua khung SockJS).
Tin nhắn gửi qua STOMP được publish kèm sự kiện Kafka \path{MESSAGE} để
notification-service báo cho người nhận đang offline. \path{MatchEventListener}
lắng nghe sự kiện \path{MATCH} từ dating-service để tự tạo \path{Conversation}
ngay khi hai người vừa ghép đôi thành công, không cần người dùng tự bắt
chuyện trước.

\section{notification-service (cổng 8093)}

\textbf{Mục đích:} tổng hợp và đẩy thông báo thời gian thực cho mọi loại sự
kiện phát sinh từ các service khác.

\subsubsection*{Cơ sở dữ liệu (MongoDB \path{notification_db})}
\begin{longtable}{p{3.2cm}p{3cm}p{8.3cm}}
\entityTableHeader{Notification}
\path{id} & \path{String} & \\
\path{recipientId} & \path{String} & Người nhận thông báo. \\
\path{actorId} & \path{String} & Người gây ra sự kiện (ví dụ người thả cảm xúc). \\
\path{type} & \path{Enum (NotificationType)} & \path{FRIEND_REQUEST}, \path{COMMENT}, \path{REACTION}, \path{GROUP}, \path{MATCH}, \path{MESSAGE}, \path{TAG}, \path{FOLLOW}. \\
\path{targetType}, \path{targetId} & \path{String} & Đối tượng liên quan (bài viết, hội thoại, ...). \\
\path{message} & \path{String} & Nội dung hiển thị. \\
\path{read} & \path{boolean} & \\
\path{createdAt} & \path{Instant} & Có \emph{TTL index} 8 tuần — MongoDB tự xoá thông báo cũ. \\
\end{longtable}

\subsubsection*{API chính}
\begin{longtable}{p{2cm}p{5.5cm}p{7cm}}
\apiTableHeader{notification-service}
\path{GET} & \path{/api/notifications/me} & Danh sách thông báo của mình (phân trang). \\
\path{PUT} & \path{/api/notifications/\{id\}/read} & Đánh dấu một thông báo đã đọc. \\
\path{PUT} & \path{/api/notifications/read-all} & Đánh dấu tất cả đã đọc. \\
\path{GET} & \path{/api/notifications/unread-count} & Số thông báo chưa đọc (hiển thị badge). \\
\path{DELETE} & \path{/api/notifications/\{id\}} & Xoá một thông báo. \\
\end{longtable}

\textbf{Ghi chú:} \path{NotificationEventListener} có 8 phương thức
\path{@KafkaListener} — lắng nghe đồng thời \path{FRIEND_REQUEST},
\path{COMMENT_CREATED}, \path{REACTION}, \path{GROUP}, \path{MATCH},
\path{MESSAGE}, \path{POST_TAGGED}, \path{FOLLOW} — biến mỗi sự kiện nghiệp vụ
từ bất kỳ service nào thành một document \path{Notification} rồi đẩy realtime
qua kênh WebSocket riêng \path{/ws-notifications} (gateway route
\path{uri: lb:ws://NOTIFICATION-SERVICE}), độc lập với kênh chat vì gateway
chỉ định tuyến được một service cho mỗi path.

\section{feed-service (cổng 8094)}

\textbf{Mục đích:} bảng tin cá nhân, dùng mô hình \emph{fan-out-on-write}
(ghi lúc đăng bài) kết hợp xếp hạng lúc đọc (ranking-on-read) để vừa nhanh vừa
không cứng nhắc theo thời gian thuần tuý.

\textbf{Lưu trữ:} không có cơ sở dữ liệu quan hệ riêng — dùng Redis \emph{sorted
set} (ZSET) cho mỗi người dùng, phần tử là ID bài viết, điểm số (score) là
thời điểm đăng (timestamp). Việc \emph{ghi} (fan-out) chỉ lưu timestamp thuần;
việc \emph{xếp hạng} chỉ tính lúc đọc — nhờ vậy đổi công thức xếp hạng không
cần migrate dữ liệu đã fan-out.

\subsubsection*{API chính}
\begin{longtable}{p{2cm}p{5.5cm}p{7cm}}
\apiTableHeader{feed-service}
\path{GET} & \path{/api/feed/me} & Bảng tin đã xếp hạng (phân trang) — xem công thức ở mục~\ref{sec:feedrank}. \\
\path{POST} & \path{/api/feed/seen} & Đánh dấu một loạt bài viết đã hiển thị (phục vụ phân trang/infinite scroll không lặp). \\
\end{longtable}

\textbf{Ghi chú:} \path{PostCreatedConsumer} lắng nghe sự kiện Kafka
\path{POST_CREATED} từ post-service, tra cứu danh sách bạn bè/người theo dõi
của tác giả (qua Feign sang user-service) rồi ghi (fan-out) ID bài viết vào
ZSET Redis của từng người nhận.

\section{moderation-service (cổng 8095)}

\textbf{Mục đích:} kiểm duyệt nội dung — vừa chủ động (lọc từ ngữ thô tục lúc
tạo nội dung), vừa bị động (người dùng báo cáo, quản trị viên duyệt và xử lý).

\subsubsection*{Cơ sở dữ liệu (\path{moderation_db})}
\begin{longtable}{p{3.2cm}p{3cm}p{8.3cm}}
\entityTableHeader{Report}
\path{id} & \path{String (UUID)} & \\
\path{reporterId} & \path{String} & Người báo cáo. \\
\path{targetType} & \path{Enum (ReportTargetType)} & Loại nội dung bị báo cáo (POST/COMMENT/REEL/STORY/GROUP/FANPAGE/USER, ...). \\
\path{targetId} & \path{String} & \\
\path{reason} & \path{Enum (ReportReason)} & \\
\path{description} & \path{Text} & Mô tả thêm của người báo cáo. \\
\path{status} & \path{Enum (ReportStatus)} & \path{PENDING} mặc định. \\
\path{reviewedBy}, \path{reviewedAt} & & Quản trị viên xử lý và thời điểm. \\
\path{resolutionNote} & \path{Text} & \\
\path{createdAt} & \path{Instant} & \\
\end{longtable}

\subsubsection*{API chính}
\begin{longtable}{p{2cm}p{5.5cm}p{7cm}}
\apiTableHeader{moderation-service}
\path{POST} & \path{/api/moderation/reports} & Người dùng gửi báo cáo vi phạm. \\
\path{GET} & \path{/api/moderation/reports/mine} & Báo cáo mình đã gửi. \\
\path{GET} & \path{/api/moderation/reports} & Hàng đợi báo cáo (chỉ \path{ADMIN}). \\
\path{GET} & \path{/api/moderation/reports/\{id\}} & Chi tiết một báo cáo (chỉ \path{ADMIN}). \\
\path{PUT} & \path{/api/moderation/reports/\{id\}/resolve} & Xử lý báo cáo (chỉ \path{ADMIN}) — nếu hành động là \path{REMOVE_CONTENT} thì phát sự kiện \path{CONTENT_REMOVED}. \\
\end{longtable}

\textbf{Ghi chú:} bộ lọc từ ngữ thô tục (tiếng Anh + tiếng Việt, giữ nguyên
dấu thanh) được nhúng trực tiếp vào bước tạo nội dung ở post-service,
comment-service, reels-service, story-service, group-service, fanpage-service
— chặn ngay tại nguồn thay vì chỉ xử lý hậu kiểm.

\section{search-service (cổng 8096)}

\textbf{Mục đích:} tìm kiếm tổng hợp trên nhiều loại đối tượng cùng lúc,
không sở hữu dữ liệu riêng mà tổng hợp qua Feign.

\subsubsection*{API chính}
\begin{longtable}{p{2cm}p{5.5cm}p{7cm}}
\apiTableHeader{search-service}
\path{GET} & \path{/api/search?q=\&limit=} & Gộp kết quả từ user-service, group-service, fanpage-service, post-service. \\
\end{longtable}

\textbf{Ghi chú:} mỗi Feign client gọi sang 4 service trên có circuit breaker
+ fallback trả danh sách rỗng riêng — một service phía sau gặp sự cố không
làm sập toàn bộ chức năng tìm kiếm, chỉ thiếu kết quả thuộc loại đó.

% ============================================================
\chapter{Các thuật toán và cơ chế nâng cao}
% ============================================================

\section{Xếp hạng bảng tin theo mức độ tương tác}
\label{sec:feedrank}

Mặc định, một bảng tin chỉ sắp theo thời gian đăng sẽ khiến nội dung tương tác
cao nhanh chóng bị chìm. \path{feed-service} tách rời việc \emph{ghi} và việc
\emph{xếp hạng}:

\begin{enumerate}[nosep]
    \item \textbf{Ghi (fan-out-on-write):} khi có bài viết mới
    (\path{POST_CREATED}), \path{FeedFanoutService} chỉ ghi cặp (ID bài viết,
    thời điểm đăng) vào Redis ZSET của từng người theo dõi/bạn bè tác giả —
    không tính điểm ở bước này.
    \item \textbf{Đọc (ranking-on-read):} \path{FeedQueryService} lấy một tập
    ứng viên giới hạn từ ZSET (theo thời gian gần nhất), sau đó chấm điểm mỗi
    bài viết theo công thức kiểu Hacker News:
\end{enumerate}

\begin{quote}
\centering
\path{score = recencyScore x (1 + log(1 + reactionCount + 2 x commentCount))}
\end{quote}

trong đó \path{recencyScore} giảm dần theo thời gian kể từ lúc đăng, còn phần
logarit khuếch đại nhẹ các bài có tương tác cao (bình luận được tính trọng số
gấp đôi cảm xúc). Danh sách ứng viên sau khi chấm điểm được sắp giảm dần rồi
mới phân trang trong bộ nhớ. Vì việc ghi không phụ thuộc công thức, đổi trọng
số/công thức xếp hạng chỉ cần sửa \path{FeedQueryService}, không cần chạy lại
migrate dữ liệu đã fan-out.

\section{Thuật toán ghép đôi (matching) cho dating}
\label{sec:matching}

\path{DatingProfile} tách bạch hai khái niệm dễ nhầm: \path{gender} (giới
tính của chính chủ hồ sơ) và \path{genderPreference} (giới tính đang tìm
kiếm). \path{GET /api/dating/candidates} thực hiện hai bước:

\begin{enumerate}[nosep]
    \item \textbf{Lọc tương thích hai chiều:} chỉ giữ lại ứng viên mà giới
    tính của họ khớp \path{genderPreference} của người xem \emph{và ngược
    lại} (giới tính người xem cũng khớp \path{genderPreference} của ứng viên).
    \item \textbf{Chấm điểm 0--100} cho các ứng viên còn lại, gồm hai phần:
\end{enumerate}

\begin{longtable}{p{2.5cm}p{2cm}p{9.5cm}}
\threecolTableHeader{Bảng chấm điểm matching}{Tiêu chí}{Điểm tối đa}{Cách tính}
Độ tuổi phù hợp hai chiều & 60 điểm & Cả hai người đều nằm trong khoảng
\path{[minAgePreference, maxAgePreference]} của đối phương thì đạt điểm tối
đa; càng lệch khỏi khoảng mong muốn của một trong hai bên, điểm càng giảm. \\
Sở thích chung & 40 điểm & Tính theo \emph{hệ số tương đồng Jaccard}
(\path{|A ∩ B| / |A ∪ B|}) trên tập \path{interests} của hai hồ sơ, nhân với
40. \\
\end{longtable}

Kết quả trả về qua \path{CandidateResponse} kèm luôn điểm số thành phần —
minh bạch thuật toán cho phía client thay vì chỉ trả một danh sách đã sắp sẵn
không rõ căn cứ. Khi hai người cùng vuốt \path{LIKE} lẫn nhau
(\path{POST /api/dating/swipe}), hệ thống tạo bản ghi \path{Match} và phát sự
kiện Kafka \path{MATCH} — được cả notification-service (báo tin) lẫn
chat-service (tự tạo cuộc trò chuyện) lắng nghe.

\section{Circuit breaker (Resilience4j) trên các Feign client}

Mọi lệnh gọi liên service qua OpenFeign đều được bọc trong một circuit
breaker Resilience4j với \path{fallbackFactory} riêng, tuân theo nguyên tắc:

\begin{itemize}[nosep]
    \item \textbf{Lỗi nghiệp vụ} (ví dụ service đích trả \path{404 Not Found}
    vì tài nguyên không tồn tại) — vẫn được rethrow đúng như vậy cho tầng gọi,
    \emph{không} bị circuit breaker nuốt mất, vì đây là phản hồi hợp lệ về
    mặt nghiệp vụ.
    \item \textbf{Lỗi hạ tầng} (service đích không phản hồi, timeout, connection
    refused) — circuit breaker mở sau một ngưỡng lỗi liên tiếp, fallback trả
    về \path{ServiceUnavailableException} (503) hoặc một giá trị mặc định an
    toàn (ví dụ danh sách rỗng cho search-service) tuỳ ngữ cảnh nghiệp vụ, để
    một service sập không kéo sập dây chuyền toàn hệ thống.
\end{itemize}

Các nơi áp dụng: comment-service, story-service, feed-service (Feign sang
user-service/post-service/reels-service), search-service (Feign sang cả 4
service dữ liệu), post-service (Feign sang user-service để kiểm tra bạn bè).
Trạng thái từng circuit breaker (CLOSED/OPEN/HALF\_OPEN) xem tại
\path{actuator/circuitbreakers} của service tương ứng.

\section{Pipeline kiểm duyệt nội dung}

Kiểm duyệt vận hành theo hai lớp bổ sung cho nhau:

\begin{enumerate}[nosep]
    \item \textbf{Chủ động (lúc tạo nội dung):} một bộ lọc từ ngữ thô tục
    (tiếng Anh + tiếng Việt, giữ nguyên dấu thanh để bắt cả biến thể có dấu)
    được gọi ngay trong luồng tạo nội dung của 6 service:
    post-service/comment-service/reels-service/story-service/group-service/
    fanpage-service — chặn (từ chối request) trước khi nội dung vi phạm kịp
    lưu vào database.
    \item \textbf{Bị động (hậu kiểm qua báo cáo):} người dùng báo cáo nội
    dung/tài khoản vi phạm (\path{POST /api/moderation/reports}); quản trị
    viên (role \path{ADMIN}) duyệt hàng đợi
    (\path{GET /api/moderation/reports}) và ra quyết định
    (\path{PUT /api/moderation/reports/\{id\}/resolve}). Nếu hành động là
    \path{REMOVE_CONTENT}, moderation-service phát sự kiện Kafka
    \path{CONTENT_REMOVED} — group-service và fanpage-service lắng nghe để
    xoá cascade luôn cả bảng thành viên/quản trị/người theo dõi liên quan,
    tránh để lại dữ liệu mồ côi.
\end{enumerate}

\section{Tìm kiếm tổng hợp (aggregation pattern)}

\path{search-service} không lưu trữ dữ liệu riêng mà đóng vai trò tầng tổng
hợp (aggregator): một request \path{GET /api/search} được "quạt ra" (fan-out)
thành 4 lệnh gọi Feign song song sang user-service, group-service,
fanpage-service, post-service, rồi gộp kết quả lại. Đây là một biến thể của
mẫu \emph{API Composition} thường gặp trong microservices khi cần một view
tổng hợp trên nhiều service mà không muốn tạo thêm một database vật lý mới
chỉ để phục vụ tìm kiếm.

% ============================================================
\chapter{Observability và DevOps}
\label{ch:observability}
% ============================================================

Một hệ thống nhiều service độc lập gần như không thể debug bằng cách đọc log
từng container thủ công — chương này mô tả bộ ba observability (tracing,
logging, metrics) và pipeline CI/CD đi kèm.

\section{Distributed tracing — Zipkin}

Mọi request đi qua API Gateway đều được gắn một \path{traceId} xuyên suốt mọi
service mà nó đi qua, kể cả những chặng bất đồng bộ qua Kafka (nhờ context
propagation của Micrometer Tracing). Sampling mặc định là \textbf{100\%}
(\path{TRACING_SAMPLING_PROBABILITY}, mặc định \path{1.0}) — hợp lý cho môi
trường demo/dev vì mọi request đều xem được trace ngay lập tức; khi triển
khai thật nên hạ xuống (ví dụ \path{0.1} tương đương 10\%) để giảm tải lưu trữ
Zipkin, chỉnh qua biến môi trường mà không cần build lại image. Giao diện xem
trace: \path{http://localhost:9411}.

\section{Log tập trung — Loki + Promtail + Grafana}

Promtail thu thập log của \emph{mọi} container Docker (mount trực tiếp
\path{/var/lib/docker/containers} và Docker socket), đẩy vào Loki; Grafana
đọc từ Loki qua mục Explore, lọc theo nhãn service (ví dụ
\path{\{service="auth-service"\}}). Mỗi dòng log nghiệp vụ đều được in kèm
\path{[traceId-spanId]} ở đầu dòng — cho phép từ một dòng log bất thường tra
ngược đúng trace tương ứng trên Zipkin để xem toàn cảnh request đó đi qua
những service nào.

\section{Metrics — Prometheus + Grafana dashboard}

Prometheus tự phát hiện (service discovery) các target cần scrape metrics
thông qua Eureka thay vì phải liệt kê thủ công từng service trong file cấu
hình — mỗi service mới thêm vào hệ thống tự động được scrape mà không cần sửa
Prometheus. Dashboard Grafana \textbf{Service Overview} được provision sẵn
(mở ra là thấy ngay, không cần tự tạo), hiển thị request rate, p99 latency,
error rate, JVM heap và CPU theo từng service. Grafana chạy ở chế độ anonymous
admin (\path{GF_AUTH_ANONYMOUS_ENABLED=true}) — không cần đăng nhập, phù hợp
môi trường demo/đồ án.

\section{CI/CD — GitHub Actions}

Pipeline (\path{.github/workflows/ci-cd.yml}) chạy tự động trên mỗi push/PR:

\begin{enumerate}[nosep]
    \item Build + test toàn bộ reactor Maven (20 module) trong một lần chạy
    duy nhất — tận dụng cơ chế build đa module để không phải build lại các
    module không đổi.
    \item Chạy integration test dùng \textbf{Testcontainers} — khởi động
    PostgreSQL/MongoDB/Kafka thật trong container tạm thời cho mỗi lần test,
    đảm bảo test sát với môi trường thật hơn so với mock.
    \item Build Docker image cho toàn bộ service.
    \item Khi merge vào nhánh \path{main}: push image lên GitHub Container
    Registry (GHCR).
\end{enumerate}

% ============================================================
\chapter{Triển khai hệ thống}
% ============================================================

\section{Kiến trúc triển khai bằng Docker Compose}

Toàn hệ thống (hạ tầng + 16 business service + 3 service hạ tầng Spring Cloud
+ frontend) được mô tả trong một file \path{docker-compose.yml} duy nhất, gồm
30 container. Bảng dưới tóm tắt cổng và phụ thuộc khởi động chính:

\begin{longtable}{p{4cm}p{2cm}p{7.5cm}}
\threecolTableHeader{Bảng cổng dịch vụ}{Service}{Cổng}{Ghi chú}
PostgreSQL & 5432 & Container \path{sma-postgres}, 10 database. \\
MongoDB & 27017 & Container \path{sma-mongodb}. \\
Redis & 6379 & Container \path{sma-redis}. \\
Kafka & 9092 & Chế độ KRaft, không cần Zookeeper. \\
MinIO & 9000 / 9001 & API S3 / Console quản trị. \\
Zipkin & 9411 & Giao diện xem trace. \\
Prometheus & 9090 & \\
Grafana & 3000 & Anonymous admin. \\
Loki & 3100 & Không expose trực tiếp cho người dùng cuối. \\
Eureka Server & 8761 & Dashboard service discovery. \\
Config Server & 8888 & \\
API Gateway & 8080 & Điểm vào duy nhất cho toàn bộ API + WebSocket. \\
auth-service & 8081 & \\
user-service & 8082 & \\
media-service & 8083 & \\
post-service & 8084 & \\
comment-service & 8085 & \\
reaction-service & 8086 & \\
story-service & 8087 & \\
reels-service & 8088 & \\
group-service & 8089 & \\
fanpage-service & 8090 & \\
dating-service & 8091 & \\
chat-service & 8092 & \\
notification-service & 8093 & \\
feed-service & 8094 & \\
moderation-service & 8095 & \\
search-service & 8096 & \\
Frontend (Flutter Web) & 3001 & Nginx phục vụ file tĩnh, forward tới cổng 80 trong container. \\
\end{longtable}

\section{Quy trình build và chạy}

Vì Dockerfile của mỗi service chỉ đơn giản \path{COPY target/*.jar app.jar}
(không build Maven bên trong image), file \path{.jar} phải được build sẵn
\emph{trước} khi build Docker image:

\begin{enumerate}[nosep]
    \item \textbf{Build toàn bộ reactor Maven:}
    \begin{quote}\path{mvn -q -B clean package -DskipTests}\end{quote}
    (chạy ở thư mục gốc dự án — build tuần tự cả 20 module theo đúng thứ tự
    phụ thuộc nhờ \path{<modules>} khai báo trong \path{pom.xml} gốc).
    \item \textbf{Build image và khởi động toàn bộ hệ thống:}
    \begin{quote}\path{docker compose up -d --build}\end{quote}
    \item Chờ khoảng 1--2 phút để PostgreSQL/MongoDB/Kafka/MinIO/Eureka khởi
    động xong trước khi các business service đăng ký thành công vào Eureka —
    kiểm tra tại \path{http://localhost:8761}. Nếu gọi API ngay trong lúc
    service đích chưa đăng ký xong, gateway trả về \path{503 Service
    Unavailable}.
\end{enumerate}

Trong quá trình phát triển, có thể chỉ chạy hạ tầng bằng Docker Compose
(\path{docker compose up postgres mongodb redis kafka minio -d}) rồi chạy
từng service bằng \path{mvn spring-boot:run} ngay trên máy, theo thứ tự:
\path{eureka-server} $\to$ \path{config-server} $\to$ các business service
$\to$ \path{api-gateway}.

\section{Cấu hình môi trường}

Các tham số vận hành quan trọng (địa chỉ Eureka/Config Server, khoá bí mật
JWT, ngưỡng rate limit, allowlist quản trị viên, tỉ lệ sampling tracing...)
đều được truyền qua biến môi trường thay vì hard-code trong image — cho phép
đổi cấu hình theo môi trường (dev/staging/production) mà không cần build lại.
Danh sách đầy đủ: xem Phụ lục~\ref{ap:env}.

% ============================================================
\chapter{Giao diện người dùng (Frontend)}
% ============================================================

Frontend được xây dựng bằng \textbf{Flutter Web} (module \path{frontend/},
build trong Docker bằng image \path{ghcr.io/cirruslabs/flutter:stable} vì máy
build không nhất thiết phải cài sẵn Flutter SDK, kết quả build tĩnh được phục
vụ qua Nginx ở cổng 3001), gọi thẳng vào API Gateway thật — không dùng dữ liệu
giả (mock).

Bố cục dọc kiểu Facebook, đã hiện thực đủ giao diện cho toàn bộ 16 business
service:

\begin{itemize}[nosep]
    \item Đăng ký / đăng nhập.
    \item Bảng tin (feed), đăng bài (kèm chia sẻ/gắn thẻ/quyền riêng tư tuỳ
    biến), bình luận (kèm trả lời), thả cảm xúc.
    \item Trang cá nhân (profile), kết bạn, theo dõi.
    \item Tìm kiếm tổng hợp, có trang kết quả riêng.
    \item Group, Fanpage.
    \item Story, Reels.
    \item Dating (vuốt/ghép đôi).
    \item Chat thời gian thực (STOMP/WebSocket), Notification thời gian thực
    (STOMP/WebSocket).
    \item Trang quản trị Moderation (chỉ hiển thị khi tài khoản có role
    \path{ADMIN}).
\end{itemize}

Frontend chưa được kiểm thử bằng trình duyệt thật và chưa kiểm thử round-trip
WebSocket thật ở thời điểm viết tài liệu này — xem thêm chi tiết trạng thái
từng phần trong \path{TODO.md} ở thư mục gốc dự án.

% ============================================================
\chapter{Hạn chế và hướng phát triển}
\label{ch:limitations}
% ============================================================

\section{Đơn giản hoá có chủ đích}

Một vài điểm đơn giản hoá được ghi nhận rõ ràng trong mã nguồn (bằng comment),
không phải thiếu sót ngoài ý muốn:

\begin{itemize}[nosep]
    \item \textbf{reaction-service / comment-service} nhận \path{targetOwnerId}
    trực tiếp từ client thay vì tự tra cứu qua Feign — tránh phải duy trì 3
    Feign client (sang post-service/comment-service/reels-service) chỉ để lấy
    một trường. Đánh đổi: tin tưởng dữ liệu client gửi lên là đúng.
    \item \textbf{fanpage-service} không có khái niệm "trang riêng tư" — đây
    là quyết định thiết kế có chủ đích (đúng theo hành vi Facebook thật: fanpage
    luôn công khai), không phải thiếu sót cần bổ sung như group-service.
    \item \textbf{Frontend} gọi thẳng API thật nhưng chưa được kiểm thử bằng
    trình duyệt/thiết bị thật, đặc biệt luồng WebSocket (chat, notification).
\end{itemize}

\section{Hướng phát triển tiếp theo}

\begin{itemize}[nosep]
    \item Bổ sung tầng cache (ví dụ Redis cache-aside) cho các API đọc nhiều
    (hồ sơ người dùng, chi tiết bài viết) để giảm tải PostgreSQL.
    \item Kiểm thử round-trip WebSocket thật (chat, notification) trên nhiều
    trình duyệt/thiết bị.
    \item Viết ứng dụng di động (Flutter mobile) tái sử dụng cùng bộ API REST
    + WebSocket đã có, không cần đổi backend.
    \item Chuyển sampling tracing về mức thấp hơn (ví dụ 10\%) và bổ sung
    alerting (Grafana Alerting / Alertmanager) dựa trên các dashboard metrics
    đã có sẵn khi triển khai môi trường thật.
    \item Cân nhắc thay thế bộ lọc từ ngữ thô tục dạng danh sách tĩnh bằng mô
    hình phân loại nội dung có khả năng học (machine learning) khi khối lượng
    nội dung đủ lớn.
\end{itemize}

% ============================================================
\appendix
\chapter{Phụ lục}
% ============================================================

\section{Bảng biến môi trường chính}
\label{ap:env}

\begin{longtable}{p{4.5cm}p{4cm}p{5cm}}
\threecolTableHeader{Bảng biến môi trường}{Biến}{Mặc định}{Áp dụng cho}
\path{EUREKA_URI} & \path{http://eureka-server:8761/eureka/} & Mọi service (trừ eureka-server). \\
\path{CONFIG_SERVER_URI} & \path{http://config-server:8888} & Mọi service dùng Config Server. \\
\path{KAFKA_BOOTSTRAP_SERVERS} & \path{kafka:9092} & Mọi service có sản xuất/tiêu thụ sự kiện Kafka. \\
\path{REDIS_HOST} / \path{REDIS_PORT} & \path{redis} / \path{6379} & api-gateway, feed-service, chat-service. \\
\path{JWT_SECRET} & (đặt trong \path{docker-compose.yml}, cần đổi khi triển khai thật) & auth-service, api-gateway. \\
\path{JWT_ACCESS_EXPIRATION_MS} & \path{900000} (15 phút) & api-gateway. \\
\path{JWT_REFRESH_EXPIRATION_MS} & \path{604800000} (7 ngày) & api-gateway. \\
\path{ZIPKIN_ENDPOINT} & \path{http://zipkin:9411/api/v2/spans} & Mọi service. \\
\path{TRACING_SAMPLING_PROBABILITY} & \path{1.0} (100\%) & Mọi service — hạ xuống khi triển khai thật. \\
\path{ADMIN_EMAILS} & \path{admin@social.app} & auth-service — allowlist tự gán role \path{ADMIN} lúc đăng ký. \\
\path{RATE_LIMIT_REPLENISH} / \path{RATE_LIMIT_BURST} & \path{50} / \path{100} & api-gateway — hầu hết route nghiệp vụ. \\
\path{RATE_LIMIT_AUTH_REPLENISH} / \path{RATE_LIMIT_AUTH_BURST} & \path{1} / \path{5} & api-gateway — riêng route \path{/api/auth/**}. \\
\path{RATE_LIMIT_WS_REPLENISH} / \path{RATE_LIMIT_WS_BURST} & \path{5} / \path{10} & api-gateway — hai route WebSocket. \\
\path{MINIO_ENDPOINT}, \path{MINIO_ACCESS_KEY}, \path{MINIO_SECRET_KEY} & \path{http://minio:9000}, \path{minioadmin}, \path{minioadmin} & media-service. \\
\path{DB_HOST} / \path{DB_PORT} & \path{postgres} / \path{5432} & Mọi service dùng PostgreSQL. \\
\path{MONGO_HOST} / \path{MONGO_PORT} & \path{mongodb} / \path{27017} & Mọi service dùng MongoDB. \\
\end{longtable}

\section{Danh sách 20 module Maven}

\begin{longtable}{p{6cm}p{7.5cm}}
\twocolTableHeader{Bảng module Maven}{Module}{Vai trò}
\path{common/common-lib} & Thư viện dùng chung: DTO, enum, tiện ích bảo mật (\path{JwtTokenProvider}) dùng lại ở nhiều service. \\
\path{infra/eureka-server} & Service discovery. \\
\path{infra/config-server} & Cấu hình tập trung. \\
\path{infra/api-gateway} & Cổng vào, xác thực, định tuyến, rate limit. \\
\path{services/auth-service} & Xác thực, phát hành JWT. \\
\path{services/user-service} & Hồ sơ, quan hệ xã hội. \\
\path{services/media-service} & Lưu trữ media. \\
\path{services/post-service} & Bài viết. \\
\path{services/comment-service} & Bình luận. \\
\path{services/reaction-service} & Cảm xúc, mục đã lưu. \\
\path{services/story-service} & Story. \\
\path{services/reels-service} & Reels. \\
\path{services/group-service} & Nhóm. \\
\path{services/fanpage-service} & Trang. \\
\path{services/dating-service} & Hẹn hò. \\
\path{services/chat-service} & Nhắn tin thời gian thực. \\
\path{services/notification-service} & Thông báo thời gian thực. \\
\path{services/feed-service} & Bảng tin. \\
\path{services/moderation-service} & Kiểm duyệt nội dung. \\
\path{services/search-service} & Tìm kiếm tổng hợp. \\
\end{longtable}

\end{document}

